% !TeX root = ../papers/nc-p3.tex
\section{高重複度の局所点での半回転と全階数での両対称性}
\label{sec:finite-high-order-half-turn}
\begin{theorem}\label{thm:finite-high-order-half-turn}
四辺が有限で全中央支持4の標準有限Frobenius窓に、平面局所完全交叉の
二つの定義式がともに四次以上となる点 $p$ があるとする。
他の局所成分の有無、反対切断の階数、印付き四空間の一致を仮定せず、
窓は両対称である。特に正則なら反対階数は8以下である。
\end{theorem}
\begin{proof}[他の局所因子を含む共通八点切断]
補題\ref{lem:high-order-point-rank-nine}の最初のブロック計算は、印 $p$ の局所環だけで
各辺の四式を与える。定理\ref{thm:finite-high-order-common-octad}の記号で
\[
      xm=ny,\quad sm=q(x,y),\quad nt=r(x,y),\quad
      q=c r+\kappa[x,y],\quad c\kappa\ne0
\]
である。ここで $c,\kappa$ は辺ごとの定数である。射影同型
$[s:x:n]\mapsto[s:cx:cn]$ のグラフとの辺の交叉 $Z$ は、
対角で残る三二次式が固有に交わるため全体で長さ8である。
印ではその局所イデアルが $\operatorname{ann}(n^2)=\operatorname{ann}(m^2)$ で、
接空間は $X$ の接空間に等しい。後者はグラフの局所方程式 $y(x)-cx$ に一次項がないことによる。

各々の残りの局所因子でもこの零化イデアルの記述を示す。
$n=0$ と $m=0$ は同じ閉点を通る。実際、例えば $n=0,m\ne0$ なら $x=0$、
次に $sm=0$ が $s\ne0$ に反する。両者が単元なら $xm=ny$ と
$sm-cnt=\kappa[x,y]$ からこの局所因子全体は同型のグラフ上にある。
残る閉点は印と異なるので、両射影の閉埋込みにより $x,y$ はいずれも非零である。
同じ一次座標変換を二組に施し、一つの成分を両方で単元として
$x=(1,X),y=(1,Y)$ と自明化する。括弧の非零定数倍を $\kappa$ に吸収する。
すると $m=n$, $n(Y-X)=0$, $n(s-c t)=\kappa(Y-X)$ である。
$m,Y$ を消した後、この滑らかな四次元座標近傍で辺の局所方程式は
\[
     n^2(s-c t),\quad nt-r(x,y),\quad R_5,\quad R_6
\]
となる。後の三式では消した変数を代入する。グラフは最初の式を $s-ct$ に置き換えて切られる。
$s-ct$ が単元なら局所辺で $n^2=0$ となり $Z$ はそこにない。
非単元なら、どちらも有限交叉なので四式は正則列である。二つのKoszul分解を比較すると
最上段は $n^2$ 倍で、局所双対により $Z$ のイデアルは $\operatorname{ann}(n^2)$ となる。
ここでは $n=m$ なので左右の零化イデアルも一致する。
以上で四辺がスキームとして同じ八点部分を持つ。
\end{proof}

\begin{proof}[局所自己同型の位数と完全跡による延長]
八点上で四つの線型同定の合成が射影的に恒等であることは、三二次式の網に
一次係数の関係がないことから定理\ref{thm:high-order-octad-bisymmetry}の通りである。
従って全位置を同じ四空間で表せる。補題\ref{lem:common-net-alternating-part}より
窓は両対称部分と四形式の和であり、どちらも二位置の巡回移動で不変である。
関係六空間は相補二位置での全縮約像なので $R_0=R_2,R_1=R_3$。
共通点グラフ上の二移動は従って自己同型 $\sigma$ をなし、その二乗は
定理\ref{thm:finite-cubic-separators}により点環上の恒等である。
$Z$ 上で隣接位置が幾何学的に一致するため、この二移動は $Z$ 上で恒等である。
印の接空間が $Z$ と全局所環とで一致するので $\sigma$ は印を固定し、接空間で恒等である。
よって局所環上でも恒等である。実際、接空間の双対の任意の基底を持ち上げ、
各持上げ $v$ を $(v+\sigma^*v)/2$ で置き換えれば、固定される元がなお極大イデアルを生成し、
局所Artin環全体を代数として生成する。

印の局所長は少なくとも16で、残る局所因子の全長は4以下である。
反転した第二辺の六式を第一辺の座標で読むと、印での二移動が恒等なのでその全局所環で消える。
各評価は第一辺の隣接語であり、非零なら補題\ref{lem:finite-adjacent-word-support}より
支持長6以上が必要なので他の因子でも消える。各辺で同じことがいえ、有限完全交叉の
$(1,1)$ 消去式の空間がちょうど六次元であることから、全ての隣接辺は互いの反転である。
命題\ref{prop:finite-flip-obstructions}の交代三次交叉の消失で両対称性が従う。
最後の階数上界は定理\ref{thm:all-bisymmetric-as-closure}による。
\end{proof}
